% =====================================================================
\section{Where Every Equation Lives in Our Project}
\label{sec:project}
% =====================================================================

\begin{center}\small
\begin{tabularx}{\linewidth}{@{}l X X X@{}}
\toprule
\textbf{Equation} & \textbf{Spacecraft reproduction} (\code{paper\_reproduction.py}) & \textbf{X3} (\code{x3\_sim/}, \code{gazebo.py}) & \textbf{MATLAB toy model} (\code{toy\_model/})\\
\midrule
(1) plant & \code{rhs()}, RK4 1\,ms & Gazebo; surrogate \code{X3.deriv()} & \code{step6} \code{rhs()}\\
(3)--(4) saturation & \code{sat\_vec()}, 2.5\,N\,m & clip 0.35\,N\,m + mixer limits & \code{step3}, \code{step6}\\
(8) FTPPF & \code{ftppf()} & \code{controllers.ftppf()} & \code{lib/ftppf.m}, \code{step1}\\
(11) $\veps_1,\veps_2$ & \code{control()} & \code{update()} & \code{lib/control\_law.m}, \code{step2}\\
(12), (22) law & \code{control()} & \code{update()} + cap (B9) & \code{lib/control\_law.m}, \code{step3}, \code{step5}\\
(14) adaptation & \code{simulate()} & \code{update()} (leakage floor) & \code{step4}\\
(15) Lyapunov & \code{lyapunov()} & -- & --\\
(B7) body rate & -- & \code{update()}, ``F11'' & --\\
(B8) outer loop & -- & \code{update()} & --\\
(B9) gain cap & -- & \code{update()}, ``F15'' & \code{step8}\\
(B10) mixer & -- & \code{allocation\_matrix()}, \code{FixedController} & \code{step7}\\
\bottomrule
\end{tabularx}
\end{center}
In the X3 column, \code{update()} is the method \code{FixedController.update()} of \code{x3\_sim/controllers.py}, which
\code{gazebo.py} calls on every pose message.

\subsection{Which barrier form each file uses}

\begin{center}\small
\begin{tabularx}{\linewidth}{@{}l l X@{}}
\toprule
\textbf{File} & \textbf{Barrier} & \textbf{Role}\\
\midrule
\code{paper\_reproduction.py} & paper, symmetric, none & reference results (all three selectable)\\
\code{x3\_sim/controllers.py}, \code{gazebo.py} & paper, symmetric, none & X3 controller (ROS parameter \code{controller\_mode})\\
\code{toy\_model/lib/control\_law.m} & paper, symmetric, none & MATLAB toy model\\
\code{Drones\_Mujoco\_Pointing\_Trial.ipynb} & paper & 3-D MuJoCo attitude demos\\
\code{index.html}, \code{app.js}, \code{simulation.js} & symmetric & interactive browser demo (illustration only)\\
\code{legacy/gazebo\_v1.py} & symmetric, $\veps$ clipped at 0.9 & first Gazebo node (crashed; kept for the record)\\
\bottomrule
\end{tabularx}
\end{center}

\subsection{From the equations to a running controller: one control step on the X3}

Every pose message triggers, in this order:
\begin{enumerate}[leftmargin=*]
  \item read the stamp, position and quaternion; keep the quaternion sign continuous;
  \item block 7: body rate from $\vq_{k-1}^*\otimes\vq_k$ and the stamp difference;
  \item block 8: reference (hover or square), acceleration command, thrust $T$, desired attitude $\vq_d$, errors $\vq_e$, $\vw_e$;
  \item block 3: $\rho_q(t)$, $\rho_\omega(t)$;\quad block 4: $\veps_1=\vq_{e,v}/\rho_q$, $\veps_2=\vw_e/\rho_\omega$;
  \item block 9 (= blocks 5 and 6 with the cap): torque $\vu$, update $\that$, clip to $\pm0.35$\,N\,m;
  \item block 10: $\vF=A^{-1}[T;\vtau]$ with thrust-priority de-saturation, $\Omega_i=\sqrt{F_i/k_T}$, publish.
\end{enumerate}
This is the paper's Algorithm~1 (Remark~8), with steps 2, 3 and 6 added for the quadcopter.
